Design of integrated circuits is split in two, analog design, and
digital design.

Digital design is highly automated. The digital functions are coded in
SystemVerilog (yes, I know there are others, but don't use those),
translated into a gate level netlist, and automatically generated
layout. Not everything is push-button automation, but most is.

Analog design, however, is manual work. We draw schematic, simulation
with a mathematical model of the real world, draw the analog layout
needed for the foundries to make the circuit, verify that we drew the
schematic and layout the same, extract parasitics, simulate again, and
in the end get a GDSII file.

When we mix analog and digital designs, we have two choices, analog on
top, or digital on top.

In analog on top we take the digital IP, and do the top level layout by
hand in analog tools.

In digital on top we include the analog IPs in the SystemVerilog, and
allow the digital tools to do the layout. The digital layout is still
orchestrated by people.

Which strategy is chosen depends on the complexity of the integrated
circuit. For medium to low level of complexity, analog on top is fine.
For high complexity ICs, then digital on top is the way to go.

Below is a description of the open source digital-on-top flow. The
analog is included into GDSII at the OpenRoad stage of the flow.

The GDSII is not sufficient to integrate the analog IP. The digital
needs to know how the analog works, what capacitance is on every digital
input, the propagation delay for digital input to digital outputs , the
relation between digital outputs and clock inputs, and the possible load
on digital outputs.

The details on timing and capacitance is covered in a Liberty file. The
behavior, or function of the analog circuit must be described in a
SystemVerilog file.

But how do we describe an analog function in SystemVerilog?
SystemVerilog is simulated in an digital simulator.


{
\centering
\aicfig{media/dig_des_tikz.png}
}

\emph{Figure 1: Mixed signal design flow, where the analog design is
captured both as a schematic for ngspice and as a SystemVerilog model
for the digital simulation}

\section{Digital simulation}\label{digital-simulation}\index{digital simulation@Digital simulation}

Conceptually, the digital simulator is easy.

\begin{itemize}
\item
  The order of execution of events at the same time-step does not matter
\item
  The system is causal. Changes in the future do not affect signals in
  the past or the now
\end{itemize}

In a digital simulator there will be an event queue, see below. From
start, set the current time step equals to the next time step. Check if
there are any events scheduled for the time step. Assume that execution
of events will add new time steps. Check if there is another time step,
and repeat.

Since the digital simulator only acts when something is supposed to be
done, they are inherently fast, and can handle complex systems.


{
\centering
\aicfig{media/eventqueue.png}
}

\emph{Figure 2: Flow chart of a digital simulator - advance to the next
time step, execute the events scheduled there, and repeat until the
queue is empty}

It's a fun exercise to make a digital simulator. On my Ph.D I wanted to
model ADCs, and first I had a look at SystemC, however, I disliked C++,
so I made
\href{https://sourceforge.net/projects/systemdotnet/}{SystemDotNet}

In SystemDotNet I implemented the event queue as a hash table, so it ran
a bit faster. See below.


{
\centering
\aicfig{media/systemdotnet.png}
}

\emph{Figure 3: Event queue implemented as a hash table in SystemDotNet,
where events at 10, 20, 21 and 30 ns are looked up by time step before
being run. Source: the SystemDotNet project on SourceForge}

\subsection{Digital Simulators}\label{digital-simulators}\index{digital simulators@Digital simulators}

There are both commercial and open source tools for digital simulation.
If you've never used a digital simulator, then I'd recommend you start
with iverilog. I've made some examples at
\href{https://github.com/wulffern/dicex/tree/main/project/verilog}{dicex}.

\textbf{Commercial}

\begin{itemize}
\tightlist
\item
  \href{https://www.cadence.com/ko_KR/home/tools/system-design-and-verification/simulation-and-testbench-verification/xcelium-simulator.html}{Cadence
  Excelium}
\item
  \href{https://eda.sw.siemens.com/en-US/ic/questa/simulation/advanced-simulator/}{Siemens
  Questa}
\item
  \href{https://www.synopsys.com/verification/simulation/vcs.html}{Synopsys
  VCS}
\end{itemize}

\textbf{Open Source} -
\href{https://github.com/steveicarus/iverilog}{iverilog/vpp} -
\href{https://www.veripool.org/verilator/}{Verilator} -
\href{https://sourceforge.net/projects/systemdotnet/}{SystemDotNet}

\subsubsection{Counter}\label{counter}

Below is an example of a counter in SystemVerilog. The code can be found
at
\href{https://github.com/wulffern/dicex/tree/main/sim/verilog/counter_sv}{counter\_sv}.

In the always\_comb section we code what will become the combinatorial
logic. In the always\_ff section we code what will become our registers.

\begin{Shaded}
\begin{Highlighting}[]
\KeywordTok{module}\NormalTok{ counter}\OperatorTok{(}
               \DataTypeTok{output}\NormalTok{ logic }\OperatorTok{[}\NormalTok{WIDTH}\DecValTok{{-}1}\OperatorTok{:}\DecValTok{0}\OperatorTok{]}\NormalTok{ out}\OperatorTok{,}
               \DataTypeTok{input}\NormalTok{ logic              clk}\OperatorTok{,}
               \DataTypeTok{input}\NormalTok{ logic              reset}
               \OperatorTok{);}

   \DataTypeTok{parameter}\NormalTok{ WIDTH }\OperatorTok{=} \DecValTok{8}\OperatorTok{;}

\NormalTok{   logic }\OperatorTok{[}\NormalTok{WIDTH}\DecValTok{{-}1}\OperatorTok{:}\DecValTok{0}\OperatorTok{]}\NormalTok{                    count}\OperatorTok{;}
\NormalTok{   always\_comb }\KeywordTok{begin}
\NormalTok{      count }\OperatorTok{=}\NormalTok{ out }\OperatorTok{+} \DecValTok{1}\OperatorTok{;}
   \KeywordTok{end}

\NormalTok{   always\_ff }\OperatorTok{@(}\KeywordTok{posedge}\NormalTok{ clk }\DataTypeTok{or} \KeywordTok{posedge}\NormalTok{ reset}\OperatorTok{)} \KeywordTok{begin}
      \KeywordTok{if} \OperatorTok{(}\NormalTok{reset}\OperatorTok{)}
\NormalTok{        out }\OperatorTok{\textless{}=} \DecValTok{0}\OperatorTok{;}
      \KeywordTok{else}
\NormalTok{        out }\OperatorTok{\textless{}=}\NormalTok{ count}\OperatorTok{;}
   \KeywordTok{end}

\KeywordTok{endmodule} \CommentTok{// counter}
\end{Highlighting}
\end{Shaded}

In the context of a digital simulator, we can think through how the
event queue will look.

When the clk or reset changes from zero to 1, then schedule an event
where if the reset is 1, then out will be zero in the next time step. If
reset is 0, then out will be \texttt{count} in the next time step.

In a time-step where \texttt{out} changes, then schedule an event to set
\texttt{count} to \texttt{out} plus one. As such, each positive edge of
the clock at least 2 events must be scheduled in the register transfer
level (RTL) simulation.

For example:

\begin{Shaded}
\begin{Highlighting}[]
\ExtensionTok{Assume} \KeywordTok{\textasciigrave{}}\ExtensionTok{clk,}\NormalTok{ reset, out = 0}\KeywordTok{\textasciigrave{}}

\ExtensionTok{Assume}\NormalTok{ event with }\KeywordTok{\textasciigrave{}}\ExtensionTok{clk}\NormalTok{ = 1}\KeywordTok{\textasciigrave{}}

\ExtensionTok{0:}\NormalTok{ Set }\KeywordTok{\textasciigrave{}}\ExtensionTok{out}\NormalTok{ = count}\KeywordTok{\textasciigrave{}}\NormalTok{ in next event }\ErrorTok{(}\ExtensionTok{1}\KeywordTok{)}

\ExtensionTok{1:}\NormalTok{ Set }\KeywordTok{\textasciigrave{}}\ExtensionTok{count}\NormalTok{ = out + 1}\KeywordTok{\textasciigrave{}}\NormalTok{ using }
   \ExtensionTok{logic} \ErrorTok{(}\ExtensionTok{may}\NormalTok{ consume multiple events}\KeywordTok{)} 

\ExtensionTok{X:}\NormalTok{ no further events}
\end{Highlighting}
\end{Shaded}

When we synthesis the code below into a
\href{https://github.com/wulffern/dicex/blob/main/sim/verilog/counter_sv/counter_netlist.v}{netlist}
it's a bit harder to see how the events will be scheduled, but we can
notice that clk and reset are still inputs, and for example the clock is
connected to d-flip-flops. The image below is the synthesized netlist

It should feel intuitive that a gate-level netlist will take longer to
simulate than an RTL, there are more events.


{
\centering
\aicfig{media/sv_counter.png}
}

\emph{Figure 4: Gate level netlist of the counter after synthesis, with
clk and reset driving eight D flip-flops and the adder logic}

\section{Transient analog simulation}\label{transient-analog-simulation}\index{transient analog simulation@Transient analog simulation}

Analog simulation is different. There is no quantized time step. How
fast ``things'' happen in the circuit is entirely determined by the time
constants, change in voltage, and change in current in the system.

It is possible to have a fixed time-step in analog simulation, for
example, we say that nothing is faster than 1 fs, so we pick that as our
time step. If we wanted to simulate 1 s, however, that's at least 1e15
events, and with 1 event per microsecond on a computer it's still a
simulation time of 31 years. Not a viable solution for all analog
circuits.

Analog circuits are also non-linear, properties of resistors,
capacitors, inductors, diodes may depend on the voltage or current
across, or in, the device. Solving for all the non-linear differential
equations is tricky.

An analog simulation engine must parse spice netlist, and setup
partial/ordinary differential equations for node matrix

The nodal matrix could look like the matrix below, \(i\) are the
currents, \(v\) the voltages, and \(G\) the conductances between nodes.

\[
\begin{pmatrix}
G_{11} &G_{12}  &\cdots  &G_{1N} \\ 
G_{21} &G_{22}  &\cdots  &G_{2N} \\ 
\vdots &\vdots  &\ddots  & \vdots\\ 
G_{N1} &G_{N2}  &\cdots  &G_{NN} 
\end{pmatrix}
\begin{pmatrix}
v_1\\ 
v_2\\ 
\vdots\\ 
v_N
\end{pmatrix}=
\begin{pmatrix}
i_1\\ 
i_2\\ 
\vdots\\ 
i_N
\end{pmatrix}
\]

The simulator, and devices model the non-linear current/voltage behavior
between all nodes

as such, the \(G\)'s may be non-linear functions, and include the
\(v\)'s and \(i\)'s.

Transient analysis use numerical methods to compute time evolution

The time step is adjusted automatically, often by proprietary
algorithms, to trade accuracy and simulation speed.

The numerical methods can be forward/backward Euler, or the others
listed below.

\begin{itemize}
\tightlist
\item
  \href{https://aquaulb.github.io/book_solving_pde_mooc/solving_pde_mooc/notebooks/02_TimeIntegration/02_01_EulerMethod.html}{Euler}
\item
  \href{https://aquaulb.github.io/book_solving_pde_mooc/solving_pde_mooc/notebooks/02_TimeIntegration/02_02_RungeKutta.html}{Runge-Kutta}
\item
  \href{https://en.wikipedia.org/wiki/Crank\%E2\%80\%93Nicolson_method}{Crank-Nicolson}
\item
  Gear \citeproc{ref-gear71}{{[}1{]}}
\end{itemize}

If you wish to learn more, I would recommend starting with the original
paper on analog transient analysis.

\href{https://www2.eecs.berkeley.edu/Pubs/TechRpts/1973/ERL-m-382.pdf}{SPICE
(Simulation Program with Integrated Circuit Emphasis)} published in 1973
by Nagel and Pederson

The original paper has spawned a multitude of commercial, free and open
source simulators, some are listed below.

If you have money, then buy Cadence Spectre. If you have no money, then
start with ngspice.

\textbf{Commercial} -
\href{https://www.cadence.com/ko_KR/home/tools/custom-ic-analog-rf-design/circuit-simulation/spectre-simulation-platform.html}{Cadence
Spectre} - \href{https://eda.sw.siemens.com/en-US/ic/eldo/}{Siemens
Eldo} -
\href{https://www.synopsys.com/implementation-and-signoff/ams-simulation/primesim-hspice.html}{Synopsys
HSPICE}

\textbf{Free} - \href{http://aimspice.com}{Aimspice} -
\href{https://www.analog.com/en/design-center/design-tools-and-calculators/ltspice-simulator.html}{Analog
Devices LTspice} - \href{https://xyce.sandia.gov}{xyce}

\textbf{Open Source} - \href{http://ngspice.sourceforge.net}{ngspice}

\section{Mixed signal simulation}\label{mixed-signal-simulation}\index{mixed signal simulation@Mixed signal simulation}

It is possible to co-simulate both analog and digital functions. An
illustration is shown below.

The system will have two simulators, one analog, with transient
simulation and differential equation solver, and a digital, with event
queue.

Between the two simulators there would be analog-to-digital, and
digital-to-analog converters.

To orchestrate the time between simulators there must be a global event
and time-step control. Most often, the digital simulator will end up
waiting for the analog simulator.

The challenge with mixed-mode simulation is that if the digital circuit
becomes to large, and the digital simulation must wait for analog
solver, then the simulation would take too long.

Most of the time, it's stupid to try and simulate complex system-on-chip
with mixed-signal , full detail, simulation.

For IPs, like an ADC, co-simulation works well, and is the best way to
verify the digital and analog.

But if we can't run mixed simulation, how do we verify analog with
digital?


{
\centering
\aicfig{media/mixed_simulator.png}
}

\emph{Figure 5: Mixed signal simulator, where a digital and an analog
simulator exchange values through DAC and ADC connect modules under a
shared event and timestep control}

\section{Analog SystemVerilog
Example}\label{analog-systemverilog-example}

\subsection{TinyTapeout TT06\_SAR}\label{tinytapeout-tt06_sar}


{
\centering
\aicfig{media/tt_um_TT06_SAR_wulffern.png}
}

\emph{Figure 6: Top level schematic of the TinyTapeout
tt\_um\_TT06\_SAR\_wulffern 8-bit SAR ADC, with the SAR core, the output
capture block and the power decoupling}


{
\centering
\aicfig{media/tt06_sar_anawave.png}
}

\emph{Figure 7: ngspice transient waveforms of the SAR, showing the
sampled analog input, the clock and uio\_out{[}0{]}, and the sarp/sarn
comparator inputs settling through the bit cycling}

\subsection{SAR operation}\label{sar-operation}\index{sar operation@SAR operation}

The key idea is to model the analog behavior to sufficient detail such
that we can verify the digital code. I think it's best to have a look at
a concrete example.

\begin{itemize}
\tightlist
\item
  Analog input is sampled when clock goes low (sarp/sarn)
\item
  uio\_out{[}0{]} goes high when bit-cycling is done
\item
  Digital output (ro) changes when uio\_out{[}0{]} goes high
\end{itemize}

\begin{Shaded}
\begin{Highlighting}[]
\CommentTok{//tt06{-}sar/src/project.v}
\KeywordTok{module}\NormalTok{ tt\_um\_TT06\_SAR\_wulffern }\OperatorTok{(}
                                \DataTypeTok{input} \DataTypeTok{wire}\NormalTok{        VGND}\OperatorTok{,}
                                \DataTypeTok{input} \DataTypeTok{wire}\NormalTok{        VPWR}\OperatorTok{,}
                                \DataTypeTok{input} \DataTypeTok{wire} \OperatorTok{[}\DecValTok{7}\OperatorTok{:}\DecValTok{0}\OperatorTok{]}\NormalTok{  ui\_in}\OperatorTok{,}
                                \DataTypeTok{output} \DataTypeTok{wire} \OperatorTok{[}\DecValTok{7}\OperatorTok{:}\DecValTok{0}\OperatorTok{]}\NormalTok{ uo\_out}\OperatorTok{,}
                                \DataTypeTok{input} \DataTypeTok{wire} \OperatorTok{[}\DecValTok{7}\OperatorTok{:}\DecValTok{0}\OperatorTok{]}\NormalTok{  uio\_in}\OperatorTok{,}
                                \DataTypeTok{output} \DataTypeTok{wire} \OperatorTok{[}\DecValTok{7}\OperatorTok{:}\DecValTok{0}\OperatorTok{]}\NormalTok{ uio\_out}\OperatorTok{,}
                                \DataTypeTok{output} \DataTypeTok{wire} \OperatorTok{[}\DecValTok{7}\OperatorTok{:}\DecValTok{0}\OperatorTok{]}\NormalTok{ uio\_oe}\OperatorTok{,}
\OtherTok{\textasciigrave{}ifdef ANA\_TYPE\_REAL}
                                \DataTypeTok{input} \DataTypeTok{real}\NormalTok{        ua\_0}\OperatorTok{,}
                                \DataTypeTok{input} \DataTypeTok{real}\NormalTok{        ua\_1}\OperatorTok{,}
\OtherTok{\textasciigrave{}else}
                                \CommentTok{// analog pins}
                                \DataTypeTok{inout} \DataTypeTok{wire} \OperatorTok{[}\DecValTok{7}\OperatorTok{:}\DecValTok{0}\OperatorTok{]}\NormalTok{  ua}\OperatorTok{,} 
\OtherTok{\textasciigrave{}endif}
                                \DataTypeTok{input} \DataTypeTok{wire}\NormalTok{        ena}\OperatorTok{,}
                                \DataTypeTok{input} \DataTypeTok{wire}\NormalTok{        clk}\OperatorTok{,}
                                \DataTypeTok{input} \DataTypeTok{wire}\NormalTok{        rst\_n}
                                \OperatorTok{);}
\end{Highlighting}
\end{Shaded}

\begin{Shaded}
\begin{Highlighting}[]
\CommentTok{//tt06{-}sar/src/tb\_ana.v}
\OtherTok{\textasciigrave{}ifdef ANA\_TYPE\_REAL}
   \DataTypeTok{real}\NormalTok{        ua\_0  }\OperatorTok{=} \DecValTok{0}\OperatorTok{;}
   \DataTypeTok{real}\NormalTok{        ua\_1 }\OperatorTok{=} \DecValTok{0}\OperatorTok{;}

\OtherTok{\textasciigrave{}else}
   \DataTypeTok{tri} \OperatorTok{[}\DecValTok{7}\OperatorTok{:}\DecValTok{0}\OperatorTok{]}\NormalTok{   ua}\OperatorTok{;}
\NormalTok{   logic       uain }\OperatorTok{=} \DecValTok{0}\OperatorTok{;}
   \KeywordTok{assign}\NormalTok{ ua }\OperatorTok{=}\NormalTok{ uain}\OperatorTok{;}
\OtherTok{\textasciigrave{}endif}

\OtherTok{\textasciigrave{}ifdef ANA\_TYPE\_REAL}
   \KeywordTok{always} \BaseNTok{\#100} \KeywordTok{begin}
\NormalTok{      ua\_0 }\OperatorTok{=} \DataTypeTok{$sin}\OperatorTok{(}\DecValTok{2}\NormalTok{*}\FloatTok{3.14}\NormalTok{*}\DecValTok{1}\OperatorTok{/}\DecValTok{7750}\NormalTok{*}\DataTypeTok{$time}\OperatorTok{);}
\NormalTok{      ua\_1 }\OperatorTok{=} \OperatorTok{{-}}\DataTypeTok{$sin}\OperatorTok{(}\DecValTok{2}\NormalTok{*}\FloatTok{3.14}\NormalTok{*}\DecValTok{1}\OperatorTok{/}\DecValTok{7750}\NormalTok{*}\DataTypeTok{$time}\OperatorTok{);}
   \KeywordTok{end}
\OtherTok{\textasciigrave{}endif}

\end{Highlighting}
\end{Shaded}

\begin{Shaded}
\begin{Highlighting}[]
\CommentTok{//tt06{-}sar/src/tb\_ana.v}
\NormalTok{   tt\_um\_TT06\_SAR\_wulffern dut }\OperatorTok{(}
\NormalTok{                                .VGND}\OperatorTok{(}\NormalTok{VGND}\OperatorTok{),}
\NormalTok{                                .VPWR}\OperatorTok{(}\NormalTok{VPWR}\OperatorTok{),}
\NormalTok{                                .ui\_in}\OperatorTok{(}\NormalTok{ui\_in}\OperatorTok{),}
\NormalTok{                                .uo\_out}\OperatorTok{(}\NormalTok{uo\_out}\OperatorTok{),}
\NormalTok{                                .uio\_in}\OperatorTok{(}\NormalTok{uio\_in}\OperatorTok{),}
\NormalTok{                                .uio\_out}\OperatorTok{(}\NormalTok{uio\_out}\OperatorTok{),}
\NormalTok{                                .uio\_oe}\OperatorTok{(}\NormalTok{uio\_oe}\OperatorTok{),}
\OtherTok{\textasciigrave{}ifdef ANA\_TYPE\_REAL}
\NormalTok{                                .ua\_0}\OperatorTok{(}\NormalTok{ua\_0}\OperatorTok{),}
\NormalTok{                                .ua\_1}\OperatorTok{(}\NormalTok{ua\_1}\OperatorTok{),}
\OtherTok{\textasciigrave{}else}
\NormalTok{                                .ua}\OperatorTok{(}\NormalTok{ua}\OperatorTok{),}
\OtherTok{\textasciigrave{}endif}
\NormalTok{                                .ena}\OperatorTok{(}\NormalTok{ena}\OperatorTok{),}
\NormalTok{                                .clk}\OperatorTok{(}\NormalTok{clk}\OperatorTok{),}
\NormalTok{                                .rst\_n}\OperatorTok{(}\NormalTok{rst\_n}\OperatorTok{)}
                                \OperatorTok{);}
\end{Highlighting}
\end{Shaded}

\begin{Shaded}
\begin{Highlighting}[]
\CommentTok{\#tt06{-}sar/src/Makefile}
\DecValTok{runa:}
\ErrorTok{    }\NormalTok{iverilog  {-}g2012 {-}o my\_design {-}c tb\_ana.fl {-}DANA\_TYPE\_REAL}
\NormalTok{    vvp {-}n my\_design}

\DecValTok{rund:}
\ErrorTok{    }\NormalTok{iverilog  {-}g2012 {-}o my\_design {-}c tb\_ana.fl}
\NormalTok{    vvp {-}n my\_design}
\end{Highlighting}
\end{Shaded}

\begin{Shaded}
\begin{Highlighting}[]
   \CommentTok{//tt06{-}sar/src/project.v}
   \CommentTok{//Main SAR loop}
\NormalTok{   always\_ff }\OperatorTok{@(}\KeywordTok{posedge}\NormalTok{ clk }\DataTypeTok{or} \KeywordTok{negedge}\NormalTok{ clk}\OperatorTok{)} \KeywordTok{begin}
      \KeywordTok{if}\OperatorTok{(\textasciitilde{}}\NormalTok{ui\_in}\OperatorTok{[}\DecValTok{0}\OperatorTok{])} \KeywordTok{begin}
\NormalTok{         state }\OperatorTok{\textless{}=}\NormalTok{ OFF}\OperatorTok{;}
\NormalTok{         tmp }\OperatorTok{=} \DecValTok{0}\OperatorTok{;}
\NormalTok{         dout }\OperatorTok{=} \DecValTok{0}\OperatorTok{;}
      \KeywordTok{end}
      \KeywordTok{else} \KeywordTok{begin}
         \KeywordTok{if}\OperatorTok{(}\NormalTok{OFF}\OperatorTok{)} \KeywordTok{begin}

         \KeywordTok{end}
         \KeywordTok{else} \KeywordTok{if}\OperatorTok{(}\NormalTok{clk }\OperatorTok{==} \DecValTok{1}\OperatorTok{)} \KeywordTok{begin}
\NormalTok{            state }\OperatorTok{=}\NormalTok{ SAMPLE}\OperatorTok{;}
         \KeywordTok{end}
         \KeywordTok{else} \KeywordTok{if}\OperatorTok{(}\NormalTok{clk }\OperatorTok{==} \DecValTok{0}\OperatorTok{)} \KeywordTok{begin}
\NormalTok{            state }\OperatorTok{=}\NormalTok{ CONVERT}\OperatorTok{;}
      \OtherTok{\textasciigrave{}ifdef}\NormalTok{ ANA\_TYPE\_REAL}
\NormalTok{            smpl }\OperatorTok{=}\NormalTok{ ua\_0 }\OperatorTok{{-}}\NormalTok{ ua\_1}\OperatorTok{;}
\NormalTok{            tmp }\OperatorTok{=}\NormalTok{ smpl}\OperatorTok{;}

            \KeywordTok{for}\OperatorTok{(}\NormalTok{int i}\OperatorTok{=}\DecValTok{7}\OperatorTok{;}\NormalTok{i}\OperatorTok{\textgreater{}=}\DecValTok{0}\OperatorTok{;}\NormalTok{i}\OperatorTok{{-}{-})} \KeywordTok{begin}
               \KeywordTok{if}\OperatorTok{(}\NormalTok{tmp }\OperatorTok{\textgreater{}=} \DecValTok{0}\OperatorTok{)} \KeywordTok{begin}
\NormalTok{                  tmp }\OperatorTok{=}\NormalTok{ tmp }\OperatorTok{{-}}\NormalTok{ lsb*}\DecValTok{2}\NormalTok{**}\OperatorTok{(}\NormalTok{i}\DecValTok{{-}1}\OperatorTok{);}
                  \KeywordTok{if}\OperatorTok{(}\NormalTok{i}\OperatorTok{==}\DecValTok{7}\OperatorTok{)}
\NormalTok{                    dout}\OperatorTok{[}\NormalTok{i}\OperatorTok{]} \OperatorTok{\textless{}=} \DecValTok{0}\OperatorTok{;}
                  \KeywordTok{else}
\NormalTok{                    dout}\OperatorTok{[}\NormalTok{i}\OperatorTok{]} \OperatorTok{\textless{}=} \DecValTok{1}\OperatorTok{;}
               \KeywordTok{end}
\end{Highlighting}
\end{Shaded}

\begin{Shaded}
\begin{Highlighting}[]

               \KeywordTok{else} \KeywordTok{begin}
\NormalTok{                  tmp }\OperatorTok{=}\NormalTok{ tmp }\OperatorTok{+}\NormalTok{ lsb*}\DecValTok{2}\NormalTok{**}\OperatorTok{(}\NormalTok{i}\DecValTok{{-}1}\OperatorTok{);}
                  \KeywordTok{if}\OperatorTok{(}\NormalTok{i}\OperatorTok{==}\DecValTok{7}\OperatorTok{)}
\NormalTok{                    dout}\OperatorTok{[}\NormalTok{i}\OperatorTok{]} \OperatorTok{=} \DecValTok{1}\OperatorTok{;}
                  \KeywordTok{else}
\NormalTok{                    dout}\OperatorTok{[}\NormalTok{i}\OperatorTok{]} \OperatorTok{=} \DecValTok{0}\OperatorTok{;}
               \KeywordTok{end}
            \KeywordTok{end}
\OtherTok{\textasciigrave{}else}
            \KeywordTok{if}\OperatorTok{(}\NormalTok{tmp }\OperatorTok{==} \DecValTok{0}\OperatorTok{)} \KeywordTok{begin}
\NormalTok{               dout}\OperatorTok{[}\DecValTok{7}\OperatorTok{]} \OperatorTok{\textless{}=} \DecValTok{1}\OperatorTok{;}
\NormalTok{               tmp }\OperatorTok{\textless{}=} \DecValTok{1}\OperatorTok{;}

            \KeywordTok{end}
            \KeywordTok{else} \KeywordTok{begin}
\NormalTok{               dout}\OperatorTok{[}\DecValTok{7}\OperatorTok{]} \OperatorTok{\textless{}=} \DecValTok{0}\OperatorTok{;}
\NormalTok{               tmp  }\OperatorTok{=} \DecValTok{0}\OperatorTok{;}
            \KeywordTok{end}
\OtherTok{\textasciigrave{}endif}

         \KeywordTok{end}
\NormalTok{         state }\OperatorTok{=}\NormalTok{ next\_state}\OperatorTok{;}
      \KeywordTok{end} \CommentTok{// else: !if(\textasciitilde{}ui\_in[0])}
   \KeywordTok{end} \CommentTok{// always\_ff @ (posedge clk)}
\end{Highlighting}
\end{Shaded}

\begin{Shaded}
\begin{Highlighting}[]
   \CommentTok{//tt06{-}sar/src/project.v}
   \KeywordTok{always} \OperatorTok{@(}\KeywordTok{posedge}\NormalTok{ done}\OperatorTok{)} \KeywordTok{begin}
\NormalTok{      state }\OperatorTok{=}\NormalTok{ DONE}\OperatorTok{;}
\NormalTok{      sampled\_dout }\OperatorTok{=}\NormalTok{ dout}\OperatorTok{;}
   \KeywordTok{end}

   \KeywordTok{always} \OperatorTok{@(}\NormalTok{state}\OperatorTok{)} \KeywordTok{begin}
      \KeywordTok{if}\OperatorTok{(}\NormalTok{state }\OperatorTok{==}\NormalTok{ OFF}\OperatorTok{)}
        \BaseNTok{\#2}\NormalTok{ done }\OperatorTok{=} \DecValTok{0}\OperatorTok{;}
      \KeywordTok{else} \KeywordTok{if}\OperatorTok{(}\NormalTok{state }\OperatorTok{==}\NormalTok{ SAMPLE}\OperatorTok{)}
        \BaseNTok{\#1}\FloatTok{.6}\NormalTok{ done }\OperatorTok{=} \DecValTok{0}\OperatorTok{;}
      \KeywordTok{else} \KeywordTok{if}\OperatorTok{(}\NormalTok{state }\OperatorTok{==}\NormalTok{ CONVERT}\OperatorTok{)}
        \BaseNTok{\#115}\NormalTok{ done }\OperatorTok{=} \DecValTok{1}\OperatorTok{;}
   \KeywordTok{end}
\end{Highlighting}
\end{Shaded}


{
\centering
\aicfig{media/tt06_sar_anawave2.png}
}

\emph{Figure 8: ngspice transient of the SAR at layout corner, with the
ramping analog input, the digital output code, the clock and the
uio\_out{[}0{]} done pulses}


{
\centering
\aicfig{media/tt06_sar_digwave.png}
}

\emph{Figure 9: The same behaviour in gtkwave from the SystemVerilog
real-number model, where uo\_out{[}7:0{]} steps as ua\_0 ramps}

\section{Summary}\label{summary}

The one-page version of this chapter:

\begin{itemize}
\tightlist
\item
  Digital simulators are fast because they only compute at events;
  analog simulators are slow because they solve the whole matrix at
  every timestep
\item
  Analog SystemVerilog models the analog blocks with real-valued signals
  inside the digital simulator - the mixed-signal chip simulates at
  digital speed
\item
  The model is a contract: it must match the schematic's behavior at the
  pins, and drift between them is found in co-simulation
\item
  A SAR ADC in TinyTapeout shows the whole flow: the same testbench
  drives the Verilog model and the transistor netlist
\item
  Verification from zero supply, every time - models make that cheap
  enough to actually do
\end{itemize}

\section{Would you like to know
more?}\label{would-you-like-to-know-more}

For more information on real-number modeling I would recommend
\href{https://youtu.be/gNpPslQZT-Y}{The Evolution of Real Number
Modeling}

\phantomsection\label{refs}
\begin{CSLReferences}{0}{0}
\bibitem[\citeproctext]{ref-gear71}
\CSLLeftMargin{{[}1{]} }%
\CSLRightInline{C. Gear, {``Simultaneous numerical solution of
differential-algebraic equations,''} \emph{IEEE Transactions on Circuit
Theory}, vol. 18, no. 1, pp. 89--95, 1971, doi:
\href{https://doi.org/10.1109/TCT.1971.1083221}{10.1109/TCT.1971.1083221}.
}

\end{CSLReferences}
